% Generated draft: requires mathematical and source review.
\documentclass[11pt]{article}
\usepackage[margin=25mm]{geometry}
\usepackage{fontspec}
\setmainfont{Latin Modern Roman}
\usepackage{amsmath,amssymb,mathtools,unicode-math}
\setmathfont{Latin Modern Math}
\usepackage{enumitem,needspace,xurl,hyperref}
\setlength{\parindent}{0pt}
\setlength{\parskip}{5pt}
\setlength{\emergencystretch}{3em}
\newcommand{\N}{\mathbb{N}}
\newcommand{\Z}{\mathbb{Z}}
\newcommand{\Q}{\mathbb{Q}}
\newcommand{\R}{\mathbb{R}}
\newcommand{\C}{\mathbb{C}}
\newcommand{\F}{\mathbb{F}}
\newcommand{\A}{\mathbb{A}}
\newcommand{\PP}{\mathbb{P}}
\newcommand{\E}{\mathbb{E}}
\newcommand{\eps}{\varepsilon}
\newcommand{\dd}{\,\mathrm{d}}
\begin{document}
\section*{UnsolvedMath ID: 30003031}
\textbf{Source problem code:} catalogue entry
\subsection{Definitions and mathematical statement}
The supplied record describes a four-part research program concerning representations of quadratic forms by a lattice and its genus. Its arithmetic setting is not fully specified, so the following formulation retains the missing choices explicitly.

Let \(\Lambda\) be the lattice under consideration. Let \(\mathcal C\) denote the intended collection of quadratic forms, with coefficient domain, rank, definiteness, determinant convention, and equivalence convention still to be specified. For an admissible determinant value \(d\), set
\[
T_d=\{T\in\mathcal C:\det(T)=d\},\qquad S_d\subseteq T_d.
\]
Write \(\Lambda\text{ represents }T\) for the representation relation intended by the source, and write \(r(\Lambda,T)\) for its intended representation number. The source does not specify whether representations are embeddings or other form-preserving maps, whether primitive representations are required, or whether representations are counted individually, modulo automorphisms, or with weights. Consequently, these symbols do not yet determine a unique counting function.

The notation \(\operatorname{gen}(\Lambda)\) denotes the genus of \(\Lambda\), ordinarily defined through local isometry in a fixed arithmetic setting. The coefficient ring, local places, and any real-place conditions needed to make this definition precise are not supplied.

With these conventions understood as unresolved input data, define
\[
\begin{aligned}
R_{\Lambda}(S_d)
&=\{T\in S_d:\Lambda\text{ represents }T\},\\
R_{\operatorname{gen}(\Lambda)}^{\mathrm{every}}(S_d)
&=\{T\in S_d:\text{for every }\Lambda'\in\operatorname{gen}(\Lambda),\
\Lambda'\text{ represents }T\}.
\end{aligned}
\]
The supplied program asks for the following determinations:
\begin{enumerate}
\item Characterize \(R_{\Lambda}(S_d)\): which forms in the subsets \(S_d\) are represented by the given lattice \(\Lambda\)?
\item Characterize \(R_{\operatorname{gen}(\Lambda)}^{\mathrm{every}}(S_d)\): which forms in \(S_d\) are represented by every lattice in the genus of \(\Lambda\)? The universal quantifier over the genus is part of the supplied question.
\item Determine the number
\[
N_{\Lambda}(d)=\#\{T\in T_d:\Lambda\text{ represents }T\}.
\]
Here the meaning and finiteness of the cardinality depend on the unspecified choice between forms themselves and equivalence classes of forms.
\item Determine the average of \(r(\Lambda,T)\) over \(T\in T_d\) and its asymptotic behavior. If \(T_d\) is finite and nonempty and uniform counting is intended, the corresponding average would be
\[
A_{\Lambda}(d)=\frac{1}{\#T_d}\sum_{T\in T_d}r(\Lambda,T).
\]
This displayed expression is conditional: the source does not specify finiteness, weights, a truncation for an infinite family, an asymptotic parameter, or a limiting regime. It also does not make clear whether the requested asymptotics concern this average, the total representation count, or individual representation numbers.
\end{enumerate}
No particular family \(S_d\), admissible range of \(d\), dimension restriction, or proposed asymptotic formula is supplied. Thus this is an explicitly incomplete formulation of the recorded program, requiring those choices before it becomes a determinate mathematical problem.
\subsection{Short English statement}
For quadratic forms of determinant \(d\), determine which are represented by a given lattice, which are represented by every lattice in its genus, how many are represented, and the average and asymptotic behavior of their representation numbers. The source leaves the arithmetic setting, counting conventions, selected families, and limiting regime unspecified; these must be supplied for a precise problem.
\subsection{Sources}
\begin{itemize}
\item Local UnsolvedMath registry: supplied source record for permanent ID \(30003031\), titled \emph{Representation of Quadratic Forms by Lattices and Genera}, in the Oberwolfach Reports Open Problems collection; exact source problem code \texttt{catalogue entry}. The record gives proposed year 2016 and does not identify an individual proposer.
\item Supplied report citation: \emph{Lattices and Applications in Number Theory} (2016), Oberwolfach Reports, DOI \texttt{10.4171/owr/2016/3}, pp.\ 87--154. Supplied source URL: \url{https://doi.org/10.4171/owr/2016/3}. The background attributes the record to workshop open-problem discussions documented in Oberwolfach Reports; no exact problem page or passage is supplied.
\item Supplied literature source: \emph{Oberwolfach Report 3/2016}, \url{https://doi.org/10.4171/owr/2016/3}.
\item Literature assessment dated exactly \texttt{2026-08-21}: the supplied curation labels the problem \texttt{partially\_solved} and states that representation and averaging questions have substantial partial results, including asymptotics in restricted dimensions and quaternionic-genus settings, while the broad four-part program remains open. This is supplied dated curation derived from the OWR record, not an independently verified literature assessment. Its time-sensitive status and cited claims require independent verification.
\item The registry labels the statement \texttt{corrected\_verified}; its accompanying verification note reports only safe notation, typography, and TeX-delimiter cleanup, with hypotheses, quantifiers, and subquestions preserved. That label does not establish mathematical completeness or independent source verification.
\end{itemize}
\end{document}
